\section{The lower bound: Theorem C and the determinant of a chain of lifts}\label{sec:det-chain}

This section proves Theorem~\ref{main:thm:C} (\lean{New.Data.integral_forms_long}) along the lines
described after it in Section~\ref{sec:main}; \emph{regular families} are defined here, and exist
by Section~\ref{sec:cg-selmer}. The height bound comes from a factorisation (Theorem~\ref{det-chain:thm:factor}) of
$Z(\vartheta)/Z(1)$, up to a power of $\Nm Q$ and ratios of Hecke $L$-values, into the ratios
$Z_j(\vartheta)/Z_j(1)$ of the $m$ single lifts $W_j\to V_j$, bounded in Section~\ref{sec:ol-values};
it rests on Rallis' inner product formula for the first lift of each row of a \emph{triangle} of
lifts and on the exchange identities of Section~\ref{sec:det-exchange} for its squares.

\subsection{Set-up}\label{det-chain:sec:setup}

The data $D$ are fixed (Definition~\ref{setup:def:data}). From Section~\ref{sec:ub-objects} we take
the \emph{chosen places} (the prime divisors of $\pi_M$), the level $J_M$, the \emph{principal
forms} on a hermitian space of rank $r$ (level $J_M$, the fixed parameter;
\lean{New.Data.principalForms}) and the \emph{good forms}, the image of $e_0$ in them, which on
$Y_0$ are the target forms; and the spaces of Definition~\ref{ub-objects:def:spaces}: $Y_j$ of rank
$r+jt$, $Y_m$ definite, $W_j$ of rank $r$ with $m-j$ minus signs
(\lean{New.Data.negCount_signsW}), for which the target parameter fits the theta lift with a
character of weight $-\kappa_j$ to the definite space $V_j=I_{r+t}$
(\lean{New.Data.fitsParams_descentPair}). Twists, $m$-fold twists (\lean{New.Data.MFoldTwist}) and
the character $\mu_j^0\vartheta^j$ of the $j$-th lift, $\mu_j^0$ of weight $-\kappa_j$ and tame part
$\omega^j$, are those of Definitions~\ref{data:def:twist} and~\ref{data:def:liftchar}. For a tame
$\vartheta$ of modulus $Q=Q(\vartheta)$ put $\Nm Q=\prod_{w\in Q}q_w$, $q_w$ the norm of $w$, and let
$\Sigma_\vartheta$ be the set of the chosen places and the places of $Q$.

\begin{definition}[\lean{New.Data.singleLift}, \lean{New.Data.goodLift},
\lean{New.Data.liftDeterminant}, \lean{New.Data.chainDeterminant}]\label{det-chain:def:det}
For $1\le j\le m$ and a tame character $\vartheta$, the \emph{single lift} $\Theta_{j,\vartheta}$
sends principal forms on $W_j$ to algebraic forms on $V_j$: it is the theta lift of the descent
pair $(W_j,V_j)$ with character $\mu_j^0\vartheta$, whose Schwartz datum is the Fock polynomial of
the highest vector at infinity, the \emph{block datum} modulo $\pi_M Q$ (the indicator of the
inclusion matrix at the chosen places, a character of the determinant of the top block at the
places of $Q$) and the vacuum elsewhere. For an $m$-fold twist $\vartheta$ let $\Theta_\vartheta$
be the chain of $m$ lifts from $Y_0$ to $Y_m$ (\lean{New.Data.lowerChainAtLevel},
Section~\ref{sec:ub-objects}) on the target forms. The \emph{chain determinant} and the \emph{lift
determinants} are
\[
  Z(\vartheta)=\det\bigl(\Theta_\vartheta^*\Theta_\vartheta\bigm|\mathcal{T}_D\bigr),\qquad
  Z_j(\vartheta)=\det\bigl(\Theta_{j,\vartheta^j}^*\Theta_{j,\vartheta^j}\bigm|
  \text{good forms on }W_j\bigr).
\]
\end{definition}

For a set $S$ of finite places let $L^S(s,\nu)$ be the Hecke $L$-function of $x\mapsto\nu(x/\bar x)$
without the Euler factors at $S\cup\bar S$. For a
character $\nu$ of $\U(1)(\A)$, an integer $n$ and a tame $\vartheta$ put
\begin{equation}\label{det-chain:eq:rho}
  \rho_n(\nu,\vartheta)=\prod_{z=1}^{n}
  \frac{\abs{L^{\Sigma_\vartheta}(z,(\nu\vartheta)^{-1})}}{\abs{L^{\Sigma_1}(z,\nu^{-1})}}
\end{equation}
(\lean{New.Data.lRatioProd}), and define the \emph{quotient ratios} and the \emph{turn ratios}
\[
  q_{ij}(\vartheta)=\rho_t\bigl(\mu_j^0/\mu_i^0,\vartheta^{j-i}\bigr)\ \ (1\le i<j\le m),\qquad
  \tau_{ij}(\vartheta)=\rho_t\bigl(\mu_{i+1}^0\mu_{j+1}^0,\vartheta^{j+1}\bigr)\ \ (0\le i,j<m)
\]
(\lean{New.Data.quotientRatio}, \lean{New.Data.turnRatio}). A \emph{critical pair} $p=(\nu,z)$, $\nu$
of weight $-c$ and $1\le z\le t$, is as in Definition~\ref{der:def:field}\,(c); it is
\emph{ramified} if the tame part is non-trivial exactly at the chosen places
(\lean{New.Data.CriticalPair.IsRamified}). Its \emph{ratio} at a tame $\vartheta$ and a sign
$\epsilon=\pm1$ is $\mathrm{rat}_p(\vartheta,\epsilon)=\abs{L^{\Sigma_\vartheta}(z,(\nu\vartheta^\epsilon)^{-1})}/\abs{L^{\Sigma_1}(z,\nu^{-1})}$,
and $\mathrm{rat}_p(\vartheta)=\mathrm{rat}_p(\vartheta,1)$
(\lean{New.Data.CriticalPair.ratio}). A function $F$ of $m$-fold twists is a \emph{ramified
ratio} if $F(\vartheta)=\rho_z(\nu,\vartheta^n)$ for a ramified pair $(\nu,z)$ and a fixed
$1\le n\le m$ (\lean{New.Data.IsRamifiedRatio}); then $F\ge0$ and $F(1)=1$. The $q_{ij}$ and the
$\tau_{ij}$ are ramified ratios (\lean{New.Data.isRamifiedRatio_quotientRatio},
\lean{New.Data.isRamifiedRatio_turnRatio}).

A function $c$ of $m$-fold twists \emph{has modulus $F$} if $\abs{c(\vartheta)}^2=C\,(\Nm
Q)^{u}F(\vartheta)^2$ for all $\vartheta$, with $C\in\R$ and $u\in\Z$ independent of $\vartheta$
(\lean{New.Data.HasAbsUpToNormPower}). Moduli multiply, and if $F\ge0$, $F(1)=1$ and $c(1)\ne0$,
then $\abs{c(\vartheta)/c(1)}=(\Nm Q)^{e/2}F(\vartheta)$ for some $e\in\Z$
(\lean{New.Data.HasAbsUpToNormPower.exists_norm_div_eq}).

\subsection{Regular families of twists}\label{det-chain:sec:regular}

A \emph{family of twists} is a map $\iota\mapsto\xi_\iota$ from the set $\mathcal{I}$ of pairs
$\iota\mid\ell$ (Section~\ref{sec:setup}) to the $m$-fold twists. The \emph{budget} is
$\beta_D=f a\log^3 b/(mk)$ (\lean{New.Data.denominatorBudget}). A critical pair $p=(\nu,z)$ is a
\emph{main pair for $j$} (\lean{New.Data.CriticalPair.IsMainPair}) if $1\le j\le m$,
$m<z\le r$, the tame part of $\nu$ is $\omega^j$, and $\kappa_j+r-t\le c\le\kappa_j$ at every
infinite place; main pairs are ramified. Test forms and false lifts are defined in
Section~\ref{sec:ol-values}.

\begin{definition}[\lean{New.Data.ClearsProjection}, \lean{New.Data.HasDenominatorsLe},
\lean{New.Data.IsDenominatorBound}, \lean{New.Data.HasSmallRatios},
\lean{New.Data.IsRegularFamily}]\label{det-chain:def:regular}
Let $\iota\mid\ell$, $\vartheta$ a twist, $1\le j\le m$ and $\beta\in\R$. Let $\mathcal{Z}$ be the
image of $\Theta_{j,\vartheta}$ on the good forms and $\mathrm{pr}$ the orthogonal projection onto
it. \emph{The denominators at $(\iota,\vartheta,j)$ are at most $\beta$} if
\begin{enumerate}[label=(\roman*)]
\item some operator $\Lambda$ on $\mathcal{Z}$ with $\norm{\iota\det\Lambda}_\ell\ge
  \exp(-\dim\mathcal{Z}\cdot r\beta)$ makes $\Lambda(\mathrm{pr}\,y)$ integral at $\iota$ for every
  form $y$ on $V_j$ that is integral at $\iota$ and orthogonal to the false lifts (this is the
  denominator bound of Definition~\ref{ol-values:def:false}(b) for $Z=\mathcal Z$, without its last
  clause on pairings; \lean{New.Selmer.DenominatorLe.clearsProjection});
\item for every main pair $p=(\nu,z)$ for $j$ and all test forms $\Xi,\Xi'$ for the lift of
  $\nu\vartheta$ from $\U(1)$ to the definite space of rank $2z$, the inner product of the
  projections of $\Xi$ and $\Xi'$ to the line of that lift has $\norm{\iota(\cdot)}_\ell\le e^\beta$.
\end{enumerate}
A family $\xi$ has \emph{small ratios} if $\Nm Q(\xi_\iota)$ is prime to $\ell$ for each
$\iota\mid\ell$, and for each $1\le j\le m$, each ramified pair $p$ and each sign $\epsilon$ the family
$\iota\mapsto\mathrm{rat}_p(\xi_\iota^j,\epsilon)^{-1}$ has height at most $x_1$. It is
\emph{regular} if it has small ratios and for each $1\le j\le m$ there are reals $\gamma_\ell$,
one for each prime, with $\sum_{\ell\in S}\gamma_\ell\le\beta_D$ for every finite set $S$ of primes,
such that the denominators at $(\iota,\xi_\iota^j,j)$ are at most $\gamma_\ell$ for every
$\iota\mid\ell$.
\end{definition}

Every admissible datum $D$ has a regular family of twists (Theorem~\ref{cg-selmer:thm:regular},
\lean{New.exists_isRegularFamily}).

\subsection{The norm of one lift}\label{det-chain:sec:onelift}

\begin{fact}[Rallis' inner product formula for sharp data, as a sum over cosets;
\lean{New.SiegelWeil.rallis_inner_product_cosetTerm}]\label{det-chain:fact:rallis}
Let $S$, $S'$ be hermitian spaces over $E$ with $2\rk S<\rk S'$, and let $v$ be a real datum of the
Fock model that carries forms of type $e$ on $S$ to forms of type $e'$ on $S'$ (an edge). There is
$A\ne0$, depending only on $S$, $S'$ and the real data, such that the following holds for all
characters $\chi$, $\chi'$ of the given weights. Let the finite datum $F$ be sharp for a level $J$ at a
set $S_0$ of places, of mass $c$: $F$ lies in the finite part, $v\otimes F$ is clean (a condition on
the support) and the vacuum off $S_0$, the matrix coefficient
$\pair{F,\omega(g)F}$ at finite points $g$ integral off $S_0$ is $\pair{F,F}\ne0$ on $J\cap K_0$ and $0$
elsewhere, and $c=\pair{F,F}/[K_0:J\cap K_0]$ (\lean{New.SiegelWeil.IsSharp}). Let $L$ be the theta lift by $v\otimes F$. Then $L$ is injective on the forms of
type $e$ fixed by $J$, and for such a form $y$ and every form $y'$ of type $e$,
$\pair{Ly',Ly}=\sum_\gamma A\,c\,\mathrm{mc}(\gamma)\pair{y',R(\gamma)y}$, an absolutely
convergent sum over the cosets $\gamma$ of the finite points modulo those integral off $S_0$, with
$\mathrm{mc}(\gamma)$ the matrix coefficient of the vacuum of the Weil representation. (Sources
that the Lean file names. For the formula: Rallis, $L$-functions and the oscillator
representation, 1987; J.-S. Li, Non-vanishing theorems for the cohomology of certain arithmetic
quotients, 1992. For the injectivity: Rallis, Injectivity properties of liftings associated to
Weil representations, 1984; J.-S. Li, 1992.)
\end{fact}

For $1\le j\le m$ and a tame $\vartheta$ let $R_{j,\vartheta}$ (\lean{New.Data.rallisOperator}) be
the operator $y\mapsto\sum_\gamma\mathrm{mc}(\gamma)\,\gamma y$ on principal forms on a space of
rank $r$, with the rank $r+t$ of $V_j$, the character of $V_j$ in the $j$-th pair and the cosets
off $\pi_MQ$; it is a spherical Hecke operator and preserves the good forms
(Section~\ref{sec:sp-moment}, \lean{New.Data.rallisOnGood}). Put
\[
  c_Q=(\Nm Q)^{rt}\prod_{w\in Q}\ \prod_{i=t+1}^{t+r}\bigl(1-q_w^{-i}\bigr)
  \qquad(\lean{New.DescentPair.twistFactor}).
\]

\begin{proposition}[Norm of one lift; \lean{New.Data.singleLift_injective},
\lean{New.Data.coe_adjoint_singleLift}, \lean{New.Data.normalizedDescent_eq_zero_iff},
\lean{New.Data.exists_adjoint_singleLift_comp_eq}]
\label{det-chain:prop:onelift}
Let $1\le j\le m$. For every tame character $\vartheta$ the lift $\Theta_{j,\vartheta}$ is
injective, and its adjoint sends a form $\Xi$ on $V_j$ to the descent
$g\mapsto\pair{\theta_W(g,\cdot),\Xi}$, the class sum of $\Xi$ against the theta kernel; so the
descent of $\Xi$ vanishes exactly when $\Xi$ is orthogonal to all lifts. There is $c_j\ne0$,
independent of $\vartheta$, with
$\Theta_{j,\vartheta}^*\Theta_{j,\vartheta}=c_j\,c_Q\,R_{j,\vartheta}$. Consequently
\begin{equation}\label{det-chain:eq:Zj}
  Z_j(\vartheta)=(c_j\,c_Q)^{\dT}\,\det\bigl(R_{j,\vartheta^j}\bigm|\mathcal{T}_D\bigr)
  \qquad(\lean{New.Data.exists_normDet_goodLift_sq_eq}).
\end{equation}
\end{proposition}

\begin{proof}
Here $\rk V_j=6r>2\rk W_j$. The block datum modulo $\pi_MQ$ is sharp for $J_M$
(\lean{New.DescentPair.isSharpFn_blockFn}): its squared norm counts one matrix at each chosen place
and, at $w\in Q$, the matrices whose top block is invertible. Fact~\ref{det-chain:fact:rallis}
gives the injectivity and the Gram form; the class sum product on $V_j$ is the index of the level
times the dimension of the weight times the Petersson product, and the index times the mass is
(\lean{New.DescentPair.gramFactor_div_card}) a factor of the chosen places, independent of
$\vartheta$, times $c_Q$. The adjoint is computed by writing the lift as the integral
$\int y(g)\,\theta_W(g,\cdot)\,dg$ and exchanging the integral with the class sum
(\lean{New.DescentPair.inner_thetaLift}). For \eqref{det-chain:eq:Zj}: the powers of $\vartheta$
have the same modulus, hence the same $c_Q$; the good forms on $W_j$ and on $Y_0$ have the same
dimension $\dT$, and $R_{j,\vartheta}$ has the same determinant on both, since every Hecke
eigensystem has the same multiplicity in both (Section~\ref{sec:sp-moment},
\lean{New.Data.targetDim_eq_finrank_goodForms_spaceW}, \lean{New.Data.det_rallisOnGood_eq}).
\end{proof}

\subsection{Triangles of lifts}\label{det-chain:sec:triangle}

The chain $\Theta_\vartheta$ leaves Weil's range after its first lift ($\rk Y_{n+1}<2\rk Y_n$ for
$n\ge1$), so Fact~\ref{det-chain:fact:rallis} does not apply to it.

\begin{definition}[\lean{New.Exchange.IsTriangle}, \lean{New.Exchange.weightedChain}]
\label{det-chain:def:triangle}
Let $V$ and $F(i,n)$, $i,n\ge0$, be inner product spaces, with isometries $j_i\colon V\to F(i,0)$,
and for a parameter $\vartheta$ let $\mathsf a_{i,n}\colon F(i,n)\to F(i,n+1)$ and
$\mathsf d_{i,n}\colon F(i,n+1)\to F(i+1,n)$ be linear maps and $w_n\ne0$ weights. They form a
\emph{triangle of size $m$} if
\begin{enumerate}[label=(\roman*)]
\item for $i+n+1<m$ there is $\beta$ with
  $w_n\,\mathsf a_{i+1,n}\mathsf d_{i,n}=(w_{n+2}/\bar\beta)\,\mathsf d_{i,n+1}\mathsf a_{i,n+1}$ on $F(i,n+1)$ (the square at
  $(i,n)$ commutes up to $\beta$);
\item for $i<m$: $\pair{z,\mathsf a_{i,0}j_iy}=\pair{\mathsf d_{i,0}z,j_{i+1}y}$, and $\mathsf d_{i,0}\mathsf a_{i,0}$ maps $j_i(V)$
  injectively into $j_{i+1}(V)$.
\end{enumerate}
The \emph{weighted chain} of row $i$ is $\mathsf X_i(n)=w_{n-1}\mathsf a_{i,n-1}\circ\dots\circ w_0\mathsf a_{i,0}$.
\end{definition}

\begin{lemma}[\lean{New.Exchange.IsTriangle.eq_zero_of_weightedChain_eq_zero},
\lean{New.Exchange.IsTriangle.exists_inner_weightedChain_eq}]\label{det-chain:lem:triangle}
In a triangle of size $m$:
\begin{enumerate}[label=(\alph*)]
\item for $i+n\le m$, $\mathsf X_i(n)$ is injective on $j_i(V)$; if $V\ne0$, all $\beta$ in (i) are
  non-zero;
\item suppose that the $\beta$ of the square at $(i,n)$ has modulus $L(i+1,i+n+2)$, that an adjoint
  law holds in each square with constants of modulus $1$
  (\lean{New.Exchange.HasSquareModuli}), and that
  $w_0\pair{\mathsf a_{i,0}j_ix,\mathsf a_{i,0}j_ix'}=C_i\pair{x,S_ix'}$ with $C_i$ of modulus $1$ and operators
  $S_i$ on $V$. Then
  $\pair{\mathsf X_0(m)j_0x,\mathsf X_0(m)j_0x'}=\varpi\,(w_0w_1^m/w_m)\,\pair{x,S_0S_1\cdots S_{m-1}x'}$ with
  $\varpi$ of modulus $\prod_{0<u<v\le m}L(u,v)$.
\end{enumerate}
\end{lemma}

\begin{proof}
Composing the squares (i) along row $i$ gives
$\mathsf X_{i+1}(n)\circ \mathsf d_{i,0}\mathsf a_{i,0}=\varsigma\cdot \mathsf d_{i,n}\circ \mathsf X_i(n+1)$ on $F(i,0)$, with $\varsigma$ a multiple of
$\bar\beta^{-1}$ for the product $\beta$ of the constants
(\lean{New.Exchange.IsTriangle.exists_chainCommutes}). (a) By induction on $n$: if
$\mathsf X_i(n+1)j_iy=0$, write $\mathsf d_{i,0}\mathsf a_{i,0}j_iy=j_{i+1}y'$; then $\mathsf X_{i+1}(n)j_{i+1}y'=0$, so $y'=0$, so
$y=0$ by (ii). If some $\beta$ were $0$, then $\varsigma=0$ and the same argument kills every $y$.
(b) Composing the adjoint laws along row $i$ gives
$\pair{v,\mathsf a_{i,n}\mathsf X_i(n)j_iz}=\varpi'\pair{\mathsf d_{i,n}v,\mathsf X_{i+1}(n)j_{i+1}z}$ for $z\in V$
(\lean{New.Exchange.IsTriangle.exists_chainAdjoint}). Take $v=\mathsf X_i(n+1)j_ix$ and use the composed
squares: $\pair{\mathsf X_i(n+1)j_ix,\mathsf X_i(n+1)j_iz}$ is a scalar times
$\pair{\mathsf X_{i+1}(n)j_{i+1}u,\mathsf X_{i+1}(n)j_{i+1}z}$, where $j_{i+1}u=\mathsf d_{i,0}\mathsf a_{i,0}j_ix$. The left
vector moves down a row and loses its first lift; the right one is only copied. By induction over
the rows the last pairing is a scalar times $\pair{u,S_{i+1}\cdots S_{m-1}z}$, and by (ii)
$\pair{u,w}=\pair{\mathsf a_{i,0}j_ix,\mathsf a_{i,0}j_iw}=(C_i/w_0)\pair{x,S_iw}$. The moduli multiply.
\end{proof}

\begin{theorem}[The triangle of the chain; \lean{New.Data.exists_isTriangle}]
\label{det-chain:thm:triangle}
The following data form a triangle of size $m$ satisfying the hypotheses of
Lemma~\ref{det-chain:lem:triangle}\,(b), with $L(u,v)=q_{uv}$. $F(i,n)$ is the space of forms on a
hermitian space of rank $r+nt$, made of a stem of rank $r$ and the blocks $i+1,\dots,i+n$ of rank
$t$ (\lean{New.Data.TriangleForm}); $F(0,n)$ holds the forms on $Y_n$, and $F(i,0)$ those on $Y_0$
or on its opposite space, according to the parity of $i$. $V$ is the space of principal forms on $Y_0$ and $j_i$ the
passage to the opposite space $i$ times. $\mathsf a_{i,n}$ is the theta lift adding the block $i+n+1$, with
the character $\mu^0_{i+n+1}\vartheta^{i+n+1}$ and the indicator of the inclusion matrix modulo
$\pi_MQ$, times a power of the splitting character (\lean{New.Data.triangleLift}); $w_n$ is the
order of $\GL_{r+nt}$ modulo $\pi_MQ$ (\lean{New.Data.triangleGLCard}). $\mathsf d_{i,n}$ is a transposed lift that removes the block $i+1$
(\lean{New.Data.triangleTransposedLift}); it stands for the adjoint of the map that adds it.
Finally $C_i$ is a constant and $S_i$ is $R_{i+1,\vartheta^{i+1}}$ divided by the number of
$(r+t)\times r$ matrices modulo $\pi_MQ$.
\end{theorem}

\begin{proof}
\emph{First lifts.} $\mathsf a_{i,0}$ goes from rank $r$ to rank $r+t>2r$, and its finite datum is sharp
for the principal level of $\pi_MQ$ (\lean{New.Data.isSharp_pointDatum}).
Fact~\ref{det-chain:fact:rallis} gives its Gram form: $w_0$ times it is a constant times
$R_{i+1,\vartheta^{i+1}}$ divided by that number of matrices
(\lean{New.Data.exists_hasSum_inner_triangleFirstLift}); and its injectivity on forms of level
$J_M$. The datum of $\mathsf d_{i,0}$ is chosen so that $\mathsf d_{i,0}$ is the adjoint of $\mathsf a_{i,0}$
(\lean{New.Exchange.IsEdgeFamily.exists_inner_modLift_eq_inner_modTransposedLift}); then
$\mathsf d_{i,0}\mathsf a_{i,0}j_iy=j_{i+1}v$ with $\pair{v,x}=\pair{\mathsf a_{i,0}y,\mathsf a_{i,0}x}$, so $v$ has level $J_M$
and $v=0$ forces $y=0$. This is (ii).
\emph{Squares.} At $(i,n)$ the two composites from $F(i,n+1)$ to $F(i+1,n+1)$, ``add the last
block, remove the first'' and ``remove the first, add the last'', are compared by the exchange
identity of Section~\ref{sec:det-exchange} (\lean{New.Exchange.one_square_identity_of_isEdgeFamily}).
With single matrices as finite data, the two composites satisfy a finite square law with constant
$\eta/c(\vartheta)$, where $\eta$ is the constant of Theorem~\ref{det-exchange:thm:square} and $\abs{c(\vartheta)}^2$ is a constant times
$\prod_{z=1}^{t}\abs{L^{\Sigma_\vartheta}(z,(\nu\vartheta^{n+1})^{-1})}^2$ times a power of residue
cardinalities, $\nu=\mu^0_{i+n+2}/\mu^0_{i+1}$
(\lean{New.Data.exists_squareLaw_triangleTransposedLift}). Linear algebra over the residue ring
(\leanfile{New/Determinant/Chain/SquareLaw}) turns this law into the relation (i) for the data
defined recursively from that of $\mathsf d_{i,0}$ (\lean{New.Data.exists_triangleLift_triangleTransposedLift_comm}).
The product of the twisted $L$-values is the untwisted one times $(\Nm Q)^{t^2}q_{i+1,i+n+2}^2$
(\lean{New.Data.CriticalPair.lValueProd_mul_inv_eq}), so $\beta$ has modulus $q_{i+1,i+n+2}$. The
adjoint law is a second commutation relation, for the transposed lifts along the reversed pairs;
these are adjoint to the descents up to constants of modulus $1$ (reversal of the theta kernels,
Fourier transform of the data), and the constants of this relation have modulus $1$
(\lean{New.Data.exists_triangleLift_triangleReversedLift_comm},
\lean{New.Data.exists_triangleTransposedLift_adjoint_comm}).
\end{proof}

\subsection{The norm and the determinant of the chain}\label{det-chain:sec:chain}

\begin{theorem}[Injectivity; \lean{New.Data.lowerChainAtLevel_injective}]\label{det-chain:thm:inj}
For every $m$-fold twist $\vartheta$, the chain $\Theta_\vartheta$ is injective on the principal
forms on $Y_0$, in particular on $\mathcal{T}_D$.
\end{theorem}

\begin{proof}
The forms $\Theta_\vartheta y$ are algebraic forms on the definite space $Y_m$, and their inner
product, a sum over classes, is $[K_0:J_\vartheta(Y_m)]$ times the dimension of the weight times
the Petersson product of the chains $\mathsf X_0(m)y$ of Theorem~\ref{det-chain:thm:triangle}
(\lean{New.Data.inner_lowerChainAtLevel}). So $\Theta_\vartheta y=0$ gives $\mathsf X_0(m)y=0$, and
$y=0$ by Lemma~\ref{det-chain:lem:triangle}\,(a).
\end{proof}

\begin{theorem}[Norm of the chain; \lean{New.Data.exists_adjoint_chain_comp_eq_smul_rallisProduct}]
\label{det-chain:thm:norm}\emergencystretch=4em
There are $C>0$, $e\in\Z$ and numbers $c(\vartheta)$ with
$\abs{c(\vartheta)}=\allowbreak C(\Nm Q)^{e/2}\prod_{1\le i<j\le m}q_{ij}(\vartheta)$ such that for every
$m$-fold twist $\vartheta$, on $\mathcal{T}_D$,
\[
  \Theta_\vartheta^*\Theta_\vartheta
  =c(\vartheta)\,c_Q^{\,m}\,R_{1,\vartheta}R_{2,\vartheta^2}\cdots R_{m,\vartheta^m}.
\]
\end{theorem}

\begin{proof}
By \lean{New.Data.inner_lowerChainAtLevel} and Lemma~\ref{det-chain:lem:triangle}\,(b) with
Theorem~\ref{det-chain:thm:triangle}, $\pair{\Theta_\vartheta x,\Theta_\vartheta y}$ is
$[K_0:J_\vartheta(Y_m)]$ times the dimension of the weight times $\varpi(\vartheta)\,w_0w_1^m/w_m$ times
$\pair{x,R_{1,\vartheta}\cdots R_{m,\vartheta^m}y}$ divided by the $m$-th power of the number of
matrices modulo $\pi_MQ$ (\lean{New.Data.exists_inner_weightedChain_eq}). Counting block triangular
groups gives $[K_0:J_\vartheta(Y_m)]=C'(\Nm Q)^{u}\,\delta^m w_m/(w_0w_1^m)$, $u\in\Z$
(\lean{New.Data.exists_stdLevelIndex_lastLevel_eq}), where $\delta$ is the factor
\lean{New.DescentPair.gramFactor} of the first single lift; $\delta$ divided by the number of
matrices modulo $\pi_MQ$ is a constant times $c_Q$. The orders $w_n$ cancel. So
the constant has modulus $\prod_{i<j}q_{ij}$, and dividing by its value at $\vartheta=1$ gives the
form of $\abs{c(\vartheta)}$. (If that value is $0$, Theorem~\ref{det-chain:thm:inj} forces
$\mathcal{T}_D=0$ and there is nothing to prove.)
\end{proof}

\begin{theorem}[Factorisation of the determinant;
\lean{New.Data.exists_chainDeterminant_pos_and_div_chainDeterminant_one_eq}]
\label{det-chain:thm:factor}
There is $e\in\Z$ such that for every $m$-fold twist $\vartheta$ the numbers $Z(\vartheta)$ and
$Z_j(\vartheta)$, $1\le j\le m$, are positive and
\[
  \frac{Z(\vartheta)}{Z(1)}=\Bigl((\Nm Q)^{e/2}\prod_{1\le i<j\le m}q_{ij}(\vartheta)\Bigr)^{\dT}
  \prod_{j=1}^{m}\frac{Z_j(\vartheta)}{Z_j(1)} .
\]
\end{theorem}

\begin{proof}
Positivity follows from the injectivity of $\Theta_\vartheta$ and $\Theta_{j,\vartheta^j}$. Take
determinants in Theorem~\ref{det-chain:thm:norm}:
$Z(\vartheta)=(c(\vartheta)c_Q^m)^{\dT}\prod_j\det(R_{j,\vartheta^j}\mid\mathcal{T}_D)$. By
\eqref{det-chain:eq:Zj}, $\prod_jZ_j(\vartheta)=(c_1\cdots c_m\,c_Q^m)^{\dT}
\prod_j\det(R_{j,\vartheta^j}\mid\mathcal{T}_D)$. The Hecke operators, whose determinants are not
computed, and the factor $c_Q$ cancel: $Z(\vartheta)=(c(\vartheta)/c_1\cdots c_m)^{\dT}\prod_jZ_j(\vartheta)$
(\lean{New.Data.exists_chainDeterminant_eq_pow_mul_prod_liftDeterminant}). At $\vartheta=1$,
$\Nm Q=1$ and $q_{ij}(1)=1$; divide.
\end{proof}

\begin{remark}
The constants $\beta$ of the squares come from the one-square identity,
whose constant $\eta$ is not computed and not shown non-zero (Theorem~\ref{det-exchange:thm:square});
they are non-zero by Lemma~\ref{det-chain:lem:triangle}\,(a), which needs the space $V$ of the
triangle (Definition~\ref{det-chain:def:triangle}) to be non-zero, and only their moduli, ratios of
$L$-values, enter.
\end{remark}

\subsection{The height of the inverse ratio}\label{det-chain:sec:height}

\begin{lemma}[\lean{New.Data.HasSmallRatios.heightLe_inv_rpow_mul_prod}]\label{det-chain:lem:ratios}
Let $\xi$ have small ratios, $e\in\Z$, and $F_i$, $i\in I$, finitely many ramified ratios. Then
$\iota\mapsto\bigl((\Nm Q(\xi_\iota))^{e/2}\prod_iF_i(\xi_\iota)\bigr)^{-1}$ has height at most
$\#I\cdot t\,x_1$.
\end{lemma}

\begin{proof}
$F_i(\vartheta)=\rho_z(\nu,\vartheta^n)$ is a product of $z\le t$ ratios of the ramified pairs
$(\nu,z')$, $z'\le z$, at $\vartheta^n$; the inverse of each has height at most $x_1$. Since
$\Nm Q(\xi_\iota)$ is prime to $\ell$, $\norm{\iota((\Nm Q(\xi_\iota))^{-e/2})}_\ell=1$, so this
factor has height $0$ (take $H=1$, $\alpha=0$, $h_\ell=0$). Heights add over products.
\end{proof}

\begin{theorem}[The inverse ratio; \lean{New.Data.heightLe_chainDeterminant_one_div},
\lean{New.Data.exists_hasSmallRatios_heightLe_liftDeterminant_div}]\label{det-chain:thm:height}
There is a family $\xi$ with small ratios such that $\iota\mapsto Z(1)/Z(\xi_\iota)$ has height at
most $\dT\,(b+f a\log^4b-\hmain)$.
\end{theorem}

\begin{proof}
Take a regular family (Theorem~\ref{cg-selmer:thm:regular}). As $k\log b\cdot\beta_D=f a\log^4b/m$,
Theorem~\ref{ol-values:thm:main} (\lean{New.Data.heightLe_liftDeterminant_ratio}) says that
$\iota\mapsto Z_j(1)/Z_j(\xi_\iota)$ has height at
most $\dT(f a\log^4b/m-k\hdir)$ for each $j$. Invert Theorem~\ref{det-chain:thm:factor}
(\lean{New.Data.exists_chainDeterminant_one_div_eq}):
\[
  \frac{Z(1)}{Z(\xi_\iota)}=\Bigl((\Nm Q(\xi_\iota))^{e/2}\prod_{i<j}q_{ij}(\xi_\iota)\Bigr)^{-\dT}
  \prod_{j=1}^{m}\frac{Z_j(1)}{Z_j(\xi_\iota)} .
\]
There are at most $m^2$ pairs $i<j$, so by Lemma~\ref{det-chain:lem:ratios} the inverse of the
bracket has height at most $m^2t\,x_1\le b$ (sizes of the data), and its $\dT$-th power at most $\dT b$. The $m$ factors
of the product add up to $\dT(f a\log^4b-mk\hdir)=\dT(f a\log^4b-\hmain)$.
\end{proof}

\subsection{Auxiliary forms}\label{det-chain:sec:aux}

\begin{definition}[\lean{New.Data.IsAuxiliaryForm}]\label{det-chain:def:aux}
Let $\gamma(\vartheta)$ be numbers, one for each $m$-fold twist. A target form $y$ is
\emph{auxiliary for $\gamma$} if $\gamma(\vartheta)\Theta_\vartheta y$ is integral at every
$\iota\mid\ell$ with $\ell\nmid\Nm Q(\vartheta)$, for every $\vartheta$, and
$\norm{\Theta_1y}_\sigma^2\le\exp(2x_1^2b^2)$ for every $\sigma\in\Aut(\C)$.
\end{definition}

\begin{theorem}[Auxiliary forms; \lean{New.Data.exists_span_isAuxiliaryForm_eq_top}]
\label{det-chain:thm:aux}
There are $e\in\Z$ and numbers $\gamma(\vartheta)\ne0$ with $\gamma(1)=1$ and
$\abs{\gamma(\vartheta)}=(\Nm Q)^{e/2}\prod_{0\le i,j<m}\tau_{ij}(\vartheta)$, such that
$\mathcal{T}_D$ is spanned by forms auxiliary for $\gamma$.
\end{theorem}

\begin{proof}
\emph{Reduction.} By \lean{New.Data.exists_span_short_lifts_eq_top}, there are $c(\vartheta)\ne0$
of modulus $\prod_{i,j}\tau_{ij}$ such that the principal forms on $Y_0$ are spanned by forms $y$
for which $(c(\vartheta)/c(1))\Theta_\vartheta y$ is integral at all $\iota\mid\ell$,
$\ell\nmid\Nm Q(\vartheta)$, and the conjugates of the values of $\Theta_1y$ have squared norm at
most $\exp(x^{202}b^2)$. Put $\gamma=c/c(1)$. As in the proof of Theorem~\ref{main:thm:A},
$\#\Gbar\,e_0$ maps the principal forms onto $\mathcal{T}_D$ and keeps integrality, and
$\norm{\Theta_1\#\Gbar e_0y}_\sigma^2\le e^{B+x^{99}}$ if the conjugate values of $\Theta_1y$ have
squared norm at most $e^B$
(\lean{New.Data.goodProjTranslate_mem_integralForms},
\lean{New.Data.normSqAt_goodProjTranslate_le}); and $x^{99}+x^{202}b^2\le2x_1^2b^2$.

\emph{Spanning forms with short lifts} (\leanfile{New/Determinant/AuxiliaryForms}).
Two square grids of hermitian spaces, the $Y^i_j$ of rank $r+(i+j)t$ with $Y^0_j=Y_j$,
and the $\tilde Y^j_i$, built on $\tilde Y^0_0=Y_0\oplus E\tilde e$ with blocks of ranks $t\pm1$,
carry on their edges theta lifts and transposed lifts with single matrices modulo
$\pi_M^{\,b}Q$ as data (\leanfile{New/Determinant/AuxiliaryForms/PathData}); a path runs through
both grids and meets the diagonal once. On the \emph{first path}
$\tilde Y^m_0,\dots,\tilde Y^0_0,Y_0,\dots,Y_m$, every principal form on $Y_0$ is some
$\operatorname{res}(\varphi)$, $\varphi$ of deep level $\pi_M^{\,b}$ on the definite space $\tilde Y^m_0$:
$\operatorname{res}(\varphi)$ has the Petersson products with forms of level $\pi_MQ$ of $\varphi$
moved down by the transposed chain and restricted to $\U(Y_0)$ at a real point
(\lean{New.Data.exists_isFlatRestriction_range_eq_top}). It is onto because
deep-level forms on $\tilde Y^0_0$ detect every non-zero principal form on $Y_0$: the period of a
deep-level Poincar\'e series is close to its identity term at a real point where the parameters
are steep (\lean{New.Data.exists_eq_zero_of_orthogonal_gridLevelForms}). The pairing of a form
on $Y_m$ with $\Theta_\vartheta\operatorname{res}(\varphi)$ is a period of a path datum on the first path
(\lean{New.Data.exists_isPathDatum_integral_eq}), and $m^2$ \emph{turns}, each crossing one square
by the exchange identity or the diagonal by two seesaws, carry it to the \emph{last
path}, where all spaces are definite; the turn at $(i,j)$ divides by a constant of modulus
$\tau_{ij}$ (\lean{New.Data.exists_transfersPeriods_first_last}). There the period is a
finite sum $\Xi$ of branching pairings of algebraic values; $n_\Xi\Xi$ is integral away from
$\Nm Q$ for some $n_\Xi$ with $\log n_\Xi\le x^{99}b^2$, and for $\vartheta=1$ its conjugates are bounded
(\lean{New.Data.exists_mul_lastPathPeriod_mem_embValRing},
\lean{New.Data.norm_lastPathPeriod_le}). Pairing with delta forms on $Y_m$ reads the values of
$\Theta_\vartheta\operatorname{res}(\varphi)$ (\lean{New.Data.exists_dualPairing_eq_mul_lastPathPeriod}), and the
deep-level forms on $\tilde Y^m_0$ are spanned by short ones
(\lean{New.Data.span_isShort_eq_top_gridLevelForms}).
\end{proof}

\subsection{Proof of Theorem C}\label{det-chain:sec:proofC}

\begin{proof}[Proof of Theorem~\ref{main:thm:C} (\lean{New.Data.integral_forms_long})]
Let $y_1,\dots,y_n$, $c$ and $\delta$ be as there. We apply Theorem~\ref{setup:thm:twist} to the twisted lattice of
Proposition~\ref{setup:prop:targetlattice} on $V=\mathcal{T}_D$ (its maps are injective by
Theorem~\ref{det-chain:thm:inj}).

\emph{The twists.} Let $\xi$ be the family of Theorem~\ref{det-chain:thm:height}. Since
$\Nm Q(\xi_\iota)$ is prime to $\ell$, the twist $\xi_\iota$ is allowed at $\ell$.

\emph{The completion $z$.} Let $\gamma$ be as in Theorem~\ref{det-chain:thm:aux}. As the auxiliary
forms span $\mathcal{T}_D$, some of them, $z_1,\dots,z_{n'}$, complete $y$ to a basis
(\lean{New.exists_extend_basis_of_span_eq_top}); then $n+n'=\dT$, so $n'\le4\dT/b^2$.

\emph{The hypotheses.} (i) holds by assumption. (ii) holds with $\gamma(\vartheta)\ne0$, by
Definition~\ref{det-chain:def:aux}. (iii) holds with $b_{\mathrm{short}}=2x_1^2b^2$
(\lean{New.Data.normSqAt_targetLattice}). For (iv),
$\abs{\gamma(\xi_\iota)}^{-2}=\bigl((\Nm Q(\xi_\iota))^{e/2}\prod_{i,j}\tau_{ij}(\xi_\iota)
\bigr)^{-2}$ is the square of a family as in Lemma~\ref{det-chain:lem:ratios} with $m^2$ ramified
ratios, so it has height at most $2m^2t\,x_1\le b=:b_{\mathrm{denom}}$
(\lean{New.Data.HasSmallRatios.heightLe_inv_rpow_mul_prod_turnRatio}); and
$b_{\mathrm{vol}}=\dT(b+f a\log^4b-\hmain)$ by Theorem~\ref{det-chain:thm:height}.

\emph{The arithmetic.} Theorem~\ref{setup:thm:twist} gives
\[
  c+2\delta+n'\,(2x_1^2b^2+b)+\dT\,(b+f a\log^4b-\hmain)\ \ge\ 0 .
\]
The completion costs at most $n'(2x_1^2b^2+b)\le(4\dT/b^2)(2x_1^2b^2+b)=\dT\cdot4(2x_1^2+b^{-1})\le
\dT b$, because $b\ge x_1^{10^6}$ (\lean{New.Data.mul_completion_cost_le}). Hence
\[
  c+2\delta\ \ge\ \dT\,(\hmain-f a\log^4b-2b)\ \ge\ \dT\,(\hmain-3f a\log^4b),
\]
using $b\le f a\log^4b$.
\end{proof}
